\documentclass[10pt,a4paper]{amsart}
\usepackage{amsmath,amssymb,amsthm,graphicx}
\usepackage{verbatim}
\newenvironment{reference}{\comment}{\endcomment}
%\newenvironment{reference}{}{}
\newenvironment{slogan}{\comment}{\endcomment}
\newenvironment{history}{\comment}{\endcomment}

% For commutative diagrams we use Xy-pic
\usepackage[all]{xy}

% We use 2cell for 2-commutative diagrams.
\xyoption{2cell}
% Theorem environments.
%
\theoremstyle{plain}
\newtheorem{theorem}[subsection]{Theorem}
\newtheorem{proposition}[subsection]{Proposition}
\newtheorem{lemma}[subsection]{Lemma}

\theoremstyle{definition}
\newtheorem{definition}[subsection]{Definition}
\newtheorem{example}[subsection]{Example}
\newtheorem{exercise}[subsection]{Exercise}
\newtheorem{situation}[subsection]{Situation}

\theoremstyle{remark}
\newtheorem{remark}[subsection]{Remark}
\newtheorem{remarks}[subsection]{Remarks}

\numberwithin{equation}{subsection}

% Macros
%
\def\lim{\mathop{\mathrm{lim}}\nolimits}
\def\colim{\mathop{\mathrm{colim}}\nolimits}
\def\Spec{\mathop{\mathrm{Spec}}}
\def\Hom{\mathop{\mathrm{Hom}}\nolimits}
\def\Ext{\mathop{\mathrm{Ext}}\nolimits}
\def\SheafHom{\mathop{\mathcal{H}\!\mathit{om}}\nolimits}
\def\SheafExt{\mathop{\mathcal{E}\!\mathit{xt}}\nolimits}
\def\Sch{\mathit{Sch}}
\def\Mor{\mathop{\mathrm{Mor}}\nolimits}
\def\Ob{\mathop{\mathrm{Ob}}\nolimits}
\def\Sh{\mathop{\mathit{Sh}}\nolimits}
\def\NL{\mathop{N\!L}\nolimits}
\def\CH{\mathop{\mathrm{CH}}\nolimits}
\def\proetale{{pro\text{-}\acute{e}tale}}
\def\etale{{\acute{e}tale}}
\def\QCoh{\mathit{QCoh}}
\def\Ker{\mathop{\mathrm{Ker}}}
\def\Im{\mathop{\mathrm{Im}}}
\def\Coker{\mathop{\mathrm{Coker}}}
\def\Coim{\mathop{\mathrm{Coim}}}

% Boxtimes
%
\DeclareMathSymbol{\boxtimes}{\mathbin}{AMSa}{"02}

%
% Macros for moduli stacks/spaces
%
\def\QCohstack{\mathcal{QC}\!\mathit{oh}}
\def\Cohstack{\mathcal{C}\!\mathit{oh}}
\def\Spacesstack{\mathcal{S}\!\mathit{paces}}
\def\Quotfunctor{\mathrm{Quot}}
\def\Hilbfunctor{\mathrm{Hilb}}
\def\Curvesstack{\mathcal{C}\!\mathit{urves}}
\def\Polarizedstack{\mathcal{P}\!\mathit{olarized}}
\def\Complexesstack{\mathcal{C}\!\mathit{omplexes}}
% \Pic is the operator that assigns to X its picard group, usage \Pic(X)
% \Picardstack_{X/B} denotes the Picard stack of X over B
% \Picardfunctor_{X/B} denotes the Picard functor of X over B
\def\Pic{\mathop{\mathrm{Pic}}\nolimits}
\def\Picardstack{\mathcal{P}\!\mathit{ic}}
\def\Picardfunctor{\mathrm{Pic}}
\def\Deformationcategory{\mathcal{D}\!\mathit{ef}}
\usepackage{cleveref}
\crefname{originalreference}{original reference}{original references}
\Crefname{originalreference}{Original reference}{Original references}
\makeatletter
\newlabel{sites-definition-point}{{00Y5}{source}{}{sourcepointer-1}{}}
\newlabel{sites-definition-point@cref}{{[originalreference][0][]00Y5}{[1][0][]source}{}{}{}}
\makeatother
\crefname{originalreference}{original reference}{original references}
\Crefname{originalreference}{Original reference}{Original references}
\makeatletter
\newlabel{sites-equation-stalk}{{04EH}{source}{}{sourcepointer-2}{}}
\newlabel{sites-equation-stalk@cref}{{[originalreference][0][]04EH}{[1][0][]source}{}{}{}}
\makeatother
\crefname{originallemma}{original lemma}{original lemmas}
\Crefname{originallemma}{Original lemma}{Original lemmas}
\makeatletter
\newlabel{lemma-geometric-lift-to-usual}{{05VN}{source}{}{sourcepointer-3}{}}
\newlabel{lemma-geometric-lift-to-usual@cref}{{[originallemma][0][]05VN}{[1][0][]source}{}{}{}}
\makeatother
\crefname{originaldefinition}{original definition}{original definitions}
\Crefname{originaldefinition}{Original definition}{Original definitions}
\makeatletter
\newlabel{definition-etale-neighbourhood}{{04JV}{source}{}{sourcepointer-4}{}}
\newlabel{definition-etale-neighbourhood@cref}{{[originaldefinition][0][]04JV}{[1][0][]source}{}{}{}}
\makeatother
\crefname{originallemma}{original lemma}{original lemmas}
\Crefname{originallemma}{Original lemma}{Original lemmas}
\makeatletter
\newlabel{lemma-geometric-lift-to-cover}{{04JX}{source}{}{sourcepointer-5}{}}
\newlabel{lemma-geometric-lift-to-cover@cref}{{[originallemma][0][]04JX}{[1][0][]source}{}{}{}}
\makeatother
\crefname{originalreference}{original reference}{original references}
\Crefname{originalreference}{Original reference}{Original references}
\makeatletter
\newlabel{sites-proposition-point-limits}{{00YC}{source}{}{sourcepointer-6}{}}
\newlabel{sites-proposition-point-limits@cref}{{[originalreference][0][]00YC}{[1][0][]source}{}{}{}}
\makeatother
\title{A geometric point of an algebraic space gives a point of its etale topos with the expected stalk functor}
\author{The Stacks Project Authors}
\date{stacks-project:spaces-properties.tex:local-snapshot-20261012; GNU FDL 1.2 or later}
\begin{document}
\maketitle
\noindent\textit{Editorial note.} All delivered mathematical text is the author's own, copied verbatim from the pinned Stacks Project chapter file spaces-properties.tex; no statement or proof step was rewritten, reordered or removed.\par\medskip
\noindent\textit{Editorial note.} Span recomputation by label pairing (LANE\_BRIEF 10.2): in this chapter the author's proof environment follows the statement environment as a standalone block; the spec covers the complete environment pair (statement span, proof span) as recorded in the pool candidate.\par\medskip
\noindent\textit{Editorial note.} Occupation scan: combined claim\_key + statement\_anchor + generic-line scan over all registration files (active and withdrawn) found 13 S10-geom published cards for spaces-properties.tex (prop-occupied.json, 2026-10-14); those 13 labels are excluded from this batch by design. For the 77 free labels: 0 occupied. The registry mentions field (7557 files) and focused-seeds (645) both point to the same S10-geom cards (dual-source check), no other occupation.\par\medskip
\ifdefined\hypertarget\hypertarget{sourcepointer-1}{}\fi
\noindent\textit{Source reference.} Original reference 00Y5: sources/stacks-chow/tags.txt, source lines 1243--1243. Cross-chapter standard prerequisite cited by the author inside the counted span; it is used as a pointer and its own proof is not part of this item.\par\medskip
\ifdefined\hypertarget\hypertarget{sourcepointer-2}{}\fi
\noindent\textit{Source reference.} Original reference 04EH: sources/stacks-chow/tags.txt, source lines 5763--5763. Cross-chapter standard prerequisite cited by the author inside the counted span; it is used as a pointer and its own proof is not part of this item.\par\medskip
\ifdefined\hypertarget\hypertarget{sourcepointer-3}{}\fi
\noindent\textit{Source reference.} Original lemma 05VN: sources/stacks-chow/tags.txt, source lines 7557--7557. Cross-chapter standard prerequisite cited by the author inside the counted span; it is used as a pointer and its own proof is not part of this item.\par\medskip
\ifdefined\hypertarget\hypertarget{sourcepointer-4}{}\fi
\noindent\textit{Source reference.} Original definition 04JV: sources/stacks-chow/tags.txt, source lines 5951--5951. Cross-chapter standard prerequisite cited by the author inside the counted span; it is used as a pointer and its own proof is not part of this item.\par\medskip
\ifdefined\hypertarget\hypertarget{sourcepointer-5}{}\fi
\noindent\textit{Source reference.} Original lemma 04JX: sources/stacks-chow/tags.txt, source lines 5953--5953. Cross-chapter standard prerequisite cited by the author inside the counted span; it is used as a pointer and its own proof is not part of this item.\par\medskip
\ifdefined\hypertarget\hypertarget{sourcepointer-6}{}\fi
\noindent\textit{Source reference.} Original reference 00YC: sources/stacks-chow/tags.txt, source lines 1250--1250. Cross-chapter standard prerequisite cited by the author inside the counted span; it is used as a pointer and its own proof is not part of this item.\par\medskip
\begin{lemma}
\label{lemma-stalk-gives-point}
\begin{slogan}
A geometric point of an algebraic space gives a point of its \'etale topos.
\end{slogan}
Let $S$ be a scheme.
Let $X$ be an algebraic space over $S$.
Let $\overline{x}$ be a geometric point of $X$.
Consider the functor
$$
u : X_\etale \longrightarrow \textit{Sets}, \quad
U \longmapsto |U_{\overline{x}}|
$$
Then $u$ defines a point $p$ of the site $X_\etale$
(Sites, Definition \ref{sites-definition-point})
and its associated stalk functor $\mathcal{F} \mapsto \mathcal{F}_p$
(Sites, Equation \ref{sites-equation-stalk})
is the functor $\mathcal{F} \mapsto \mathcal{F}_{\overline{x}}$
defined above.
\end{lemma}

\begin{proof}
In the proof of
Lemma \ref{lemma-geometric-lift-to-usual}
we have seen that the scheme $U_{\overline{x}} = \overline{x} \times_X U$
is a disjoint union of
schemes isomorphic to $\overline{x}$. Thus we can also think of
$|U_{\overline{x}}|$ as the set of geometric points of $U$ lying over
$\overline{x}$, i.e., as the collection of morphisms
$\overline{u} : \overline{x} \to U$ fitting into the diagram of
Definition \ref{definition-etale-neighbourhood}.
From this it follows that $u(X)$ is a singleton, and that
$u(U \times_V W) = u(U) \times_{u(V)} u(W)$
whenever $U \to V$ and $W \to V$ are morphisms in $X_\etale$.
And, given a covering $\{U_i \to U\}_{i \in I}$ in $X_\etale$ we see
that $\coprod u(U_i) \to u(U)$ is surjective by
Lemma \ref{lemma-geometric-lift-to-cover}.
Hence
Sites, Proposition \ref{sites-proposition-point-limits}
applies, so $p$ is a point of the site $X_\etale$.
Finally, the our functor $\mathcal{F} \mapsto \mathcal{F}_{\overline{s}}$
is given by exactly the same colimit as the functor
$\mathcal{F} \mapsto \mathcal{F}_p$ associated to $p$ in
Sites, Equation \ref{sites-equation-stalk}
which proves the final assertion.
\end{proof}

\end{document}
